\section{Nuevos paradigmas de cómputo: llevar el compute hacia los data}

En la segunda mitad de la clase se vuelve a la hardware research y se discuten in-memory, near-memory, photonics y compound semiconductor. Su objetivo común no es simplemente aumentar los peak FLOPS, sino reducir el data movement, la electrical/optical conversion y la synchronization. El lector debe ver estas direcciones como un workload-specific design space, y no como la lista definitiva de la próxima arquitectura de propósito general.

\subsection{In-memory y near-memory computing}

\term{in-memory computing} (cómputo en memoria) hace que el memory device ejecute parte de las operaciones en el mismo lugar donde almacena los datos; \term{near-memory computing} (cómputo cercano a la memoria) coloca el compute en el memory stack o junto a él, para acortar la ruta de transferencia. Ambos buscan mitigar el von Neumann data movement, pero enfrentan limitaciones de precision, device variation, endurance, programming model y workload mapping.

\lecturefigure{15-compute-near-memory.png}{El cómputo cercano a la memoria o en memoria reduce el movimiento de datos y, con ello, intercambia eficiencia energética por generalidad.}{Subtítulos locales 45:10--47:20; Sebastian et al. 2020.}

\begin{knowledgebox}{Lectura de la figura: por qué «mover menos datos» no equivale a ahorrar automáticamente la mitad de la energía}
La ganancia depende de si la operación puede mapearse a una memory primitive, de si la precisión de la conversión es suficiente, de si los ADC/DAC periféricos y el control son costosos, de si los datos deben moverse entre capas y de si el compiler puede aprovechar el hardware. Los porcentajes mencionados en clase deben entenderse como motivación de investigación, no como una garantía para cualquier workload.
\end{knowledgebox}

\subsection{Términos clave: Photonic, CXL y collectives}

\begin{knowledgebox}{Tres conceptos de interconexión que se confunden con facilidad}
\begin{tabularx}{\textwidth}{L{0.19\textwidth}L{0.25\textwidth}X}
\toprule
Término & Problema que resuelve & Mecanismo y límites \\
\midrule
Photonic interconnect & Reducir el consumo energético y la presión de densidad en transferencias de gran ancho de banda a larga distancia & Transmite datos mediante luz, pero la electro-optical conversion, el laser, el packaging y el control siguen teniendo un costo. \\
CXL & Coherent attachment / pooling de CPU, accelerator y memory & Mejora la combinación y el uso compartido de recursos; no equivale a eliminar la latency ni la contention de la memory remota. \\
Collectives & Comunicación colectiva entre varias GPU: all-reduce, all-gather, reduce-scatter, broadcast, entre otras & Es la semántica de software del distributed training/serving; el rendimiento depende de la topology, el message size, el overlap y la library. \\
\bottomrule
\end{tabularx}
\end{knowledgebox}

\begin{warningbox}{La architecture religion es tan peligrosa como la model religion}
Una arquitectura nueva puede ser muy potente en un kernel específico y, aun así, carecer de compiler, debugger, library, availability o multi-tenant isolation. Tanto en la adquisición como en la investigación conviene comparar el end-to-end workload, el migration cost y la ecosystem maturity, en lugar de basarse solo en demostraciones de dispositivos o en la eficiencia energética pico.
\end{warningbox}

\teachervoice{La invitación del ponente a los estudiantes fue directa: prestar atención al punto donde convergen las señales de publication, prototype y startup. Estas notas la reformulan como método de investigación: primero identificar el cuello de botella del sistema, después elegir un workload medible y, por último, confirmar que la ganancia del dispositivo se sostiene a través del compiler y del application stack.}

\subsection{Resumen de la sección}

Los nuevos paradigmas de cómputo giran en torno al data movement. In/near-memory, photonics, CXL y los nuevos materiales resuelven problemas de distintos niveles; la validación final sigue siendo la eficiencia energética, la corrección, la programabilidad y la confiabilidad de un workload de extremo a extremo.
